\chapter{Introduction}
\label{chap:Intro}

Unilateral spatial neglect (USN) is a neuropsychological disorder in which patients fail to attend, respond, or orient toward one side of space after brain injury. The disorder is clinically important because it can affect everyday behaviour, rehabilitation progress, and spatial interaction with the surrounding environment. Prism adaptation therapy is one established approach to USN rehabilitation. In conventional prism adaptation, participants perform repeated pointing or reaching movements while wearing prism glasses that laterally displace the visual field. With repeated exposure, pointing errors are reduced, and when the prisms are removed, participants may show an aftereffect in the opposite direction of the displacement.

Virtual reality (VR) makes it possible to reproduce parts of this adaptation logic without physical prism glasses. Instead of changing the visual field through optics, a VR system can shift a virtual hand, displace a rendered scene, rotate a pointing direction, or alter the relation between physical movement and visual feedback in software. This flexibility is useful, but it also creates a methodological problem. Across the literature, different systems are described using overlapping terms such as virtual prism adaptation, visuomotor adaptation, virtual hand displacement, and visual shift, even though the underlying implementations can differ substantially. A system that shifts the entire scene is not necessarily equivalent to one that shifts only the virtual hand or rotates the pointing vector.

This project responds to that problem by developing a configurable VR sandbox for comparing several implementations of a prism effect inside one shared task and logging framework. The aim is not to prove clinical efficacy, but to investigate whether different VR-based perturbations can be implemented, tested, logged, and analyzed in a comparable structure. The prototype therefore focuses on three prism-like perturbation implementations: Translation, Rotation, and Skew. These are treated as implementation categories rather than as claims that all three manipulations are physiologically identical.

The project also compares two input modes: controller-based pointing and embodied hand-tracking-based pointing. This distinction is important because VR prism-adaptation systems may use different forms of visual feedback and body representation. Embodied input may feel closer to natural hand movement, but it can introduce tracking and posture variability. Controller input may be more robust, but confirmation and controller orientation can still affect endpoint measurements. For this reason, the prototype uses a shared confirmation and logging structure so that input mode can be examined as part of the experimental setup rather than hidden inside implementation details.

The final prototype was evaluated with healthy adult participants. The evaluation used a baseline--exposure--post structure: first, participants completed a measurement task without perturbation; next, they completed an exposure block under the active perturbation; finally, the same measurement task was repeated without perturbation. The two main assessment tasks were Open-Loop Pointing and Line Bisection. The study also logged exposure quality, hit rates, spatial response distributions, movement samples, and participant-level summaries through an RShiny analysis dashboard.

The report is structured as follows. Chapter~\ref{chap:Background} introduces USN, conventional prism adaptation, VR-based prism-adaptation research, and the terminology problem motivating the project. Chapter~\ref{chap:Meth} defines the methodological structure used to compare perturbations and outcome measures. Chapter~\ref{chap:Design} describes the participant-facing workflow and task decomposition. Chapter~\ref{chap:implementation} explains how the Unity prototype, logging system, and RShiny analysis layer were implemented. Chapter~\ref{chap:Evaluation} describes the participant evaluation procedure. Chapter~\ref{chap:results} presents the recorded results, while Chapter~\ref{chap:Discussion} discusses their interpretation, limitations, and methodological implications. Finally, Chapter~\ref{chap:conclusion-future-work} concludes the project and outlines future work.
